\documentclass[11pt]{article}
\usepackage[utf8]{inputenc}
\usepackage[T1]{fontenc}
\usepackage{lmodern}
\usepackage{amsmath,amssymb,amsthm,mathtools}
\usepackage[margin=1in]{geometry}
\usepackage{setspace}
\usepackage{parskip}
\usepackage{titlesec}
\usepackage{xcolor}
\usepackage{hyperref}
\hypersetup{colorlinks=true,linkcolor=blue,citecolor=blue,urlcolor=blue}
\onehalfspacing

\theoremstyle{plain}
\newtheorem{theorem}{Theorem}[section]
\newtheorem{proposition}[theorem]{Proposition}
\newtheorem{corollary}[theorem]{Corollary}
\theoremstyle{definition}
\newtheorem{definition}[theorem]{Definition}
\theoremstyle{remark}
\newtheorem*{remark}{Remark}
\newtheorem*{significance}{Significance}

\newcommand{\Deff}{D_{\mathrm{eff}}}
\newcommand{\Mtau}{M_{\tau}}
\newcommand{\Rmin}{R_{\min}}
\newcommand{\Dmin}{D_{\min}}
\newcommand{\Dmax}{D_{\max}}
\newcommand{\Mmin}{M_{\min}}
\newcommand{\Mmax}{M_{\max}}
\newcommand{\Gmax}{G_{\max}}
\newcommand{\Kcop}{K_c^{\mathrm{op}}}
\newcommand{\Lobs}{L_{\mathrm{obs}}}

\title{Supplementary Text S1\\Mathematical Appendix for the Gamma-Code Consensus Model}
\author{Mehmet Oktay Tugan}
\date{}

\begin{document}
\maketitle

\noindent This supplementary text contains the mathematical well-posedness results, observable bounds, finite-horizon stability estimates, sufficient selective-preservation re-entry condition, and conservativity discussion referenced in the main manuscript.

\appendix
\makeatletter
% Inline implementation of removefromreset (avoids needing the remreset package):
% remove section from the reset-list of the theorem counter so numbering is globally unique
\providecommand*{\@removefromreset}[2]{%
  \let\@tempa\@elt
  \let\@elt\relax
  \expandafter\let\expandafter\@tempb\csname cl@#2\endcsname
  \expandafter\def\csname cl@#2\endcsname{}%
  \def\@elt##1{%
    \ifx ##1#1\else
      \expandafter\@temptokena\expandafter{\csname cl@#2\endcsname}%
      \edef\x{\noexpand\expandafter\noexpand\def\noexpand\csname cl@#2\noexpand\endcsname{\the\@temptokena \noexpand\@elt{##1}}}\x
    \fi}%
  \@tempb
  \let\@elt\@tempa}
\@removefromreset{theorem}{section}
\makeatother
\section*{Mathematical Appendix}
\addcontentsline{toc}{section}{Mathematical Appendix}
\label{sec:appendix}
\setcounter{section}{0}
\setcounter{theorem}{0}
\renewcommand{\thesection}{A.\arabic{section}}
\renewcommand{\thesubsection}{A.\arabic{section}.\arabic{subsection}}
\renewcommand{\thetheorem}{A.\arabic{theorem}}

% Redefine section format for appendix: no trailing dot (already in thesection)
\titleformat{\section}{\large\bfseries}{\thesection}{0.6em}{}
\titlespacing*{\section}{0pt}{1.5em}{0.8em}

% Widen TOC number column for appendix entries (A.1 needs more space than 1)
\addtocontents{toc}{\string\setlength{\string\cftsecnumwidth}{3.2em}}

This appendix supplements the GCC with mathematical results that address three common lines of attack: (1)~Is the dynamics well-posed? (2)~Are the observables controlled? (3)~Can the gain-based re-entry idea be formulated as a sufficient condition?


\section{Assumptions}
\label{app:voraussetzungen}

We consider, over a finite horizon $t \in [0,T]$, the GCC dynamics defined in the main text under the following assumptions.

\paragraph{A1.\;Structure.}
$W_{ij}(t) \geq 0$, \;\;$\sum_j W_{ij}(t) = 1$ for all $i$ with non-empty rows, \;\;$W_{ii}(t) = 0$.

\paragraph{A2.\;Boundedness.}
There exist finite constants $\Gmax$, $\Sigma_{\max}$, $\Phi_{\max}$ such that
\[
0 \leq G(t) \leq \Gmax, \quad
0 \leq g_i(t) \leq 1, \quad
0 \leq \sigma_i(t) \leq \Sigma_{\max}, \quad
|\varphi_{ij}(t)| \leq \Phi_{\max}.
\]

\paragraph{A3.\;Temporal regularity.}
All control functions are measurable and bounded; for stability statements, piecewise continuity suffices.

\paragraph{A4.\;Observation windows.}
$\Deff^{\lambda}$ and $\Mtau$ are computed on fixed window lengths $\tau_D > 0$ and $\tau_M > 0$, with the ridge regularization defined in the main text.


\section{Elementary Bounds of the Observables}

\begin{proposition}[Elementary bounds]
\label{prop:schranken}
For all times $t$:
\begin{equation}
0 \leq R(t) \leq 1, \qquad
1 \leq \Deff(t) \leq N, \qquad
0 \leq \Mtau(t) \leq \tfrac{1}{4}.
\end{equation}
\end{proposition}

\begin{proof}[Proof sketch]
$R(t)$ is the modulus of a mean of unit complex numbers. $\Deff(t)$ is the participation ratio of a positive semi-definite matrix. $\Mtau(t)$ is the variance of a time series with values in $[0,1]$; its maximum is $\frac{1}{4}$.
\end{proof}

\begin{significance}
All three core observables are mathematically compact and controlled. This prevents empirical fits from drifting into unphysical regions.
\end{significance}


\section{Well-Posedness of the Dynamics}

\begin{theorem}[Global Lipschitz property and well-posedness]
\label{satz:lipschitz}
Define the drift vector $b(\theta,t)$ component-wise by
\begin{equation}
b_i(\theta,t) = \omega_i + K\,G(t)\,g_i(t)\sum_j W_{ij}(t)\,\sin\!\bigl(\theta_j - \theta_i - \varphi_{ij}(t)\bigr).
\end{equation}
Then, for all $\theta,\psi \in \mathbb{R}^N$ and all $t \in [0,T]$,
\begin{equation}
\label{eq:lipschitz}
\|b(\theta,t) - b(\psi,t)\|_{\infty} \;\leq\; 2\,|K|\,\Gmax\,\|\theta - \psi\|_{\infty}.
\end{equation}
Moreover, $\|b(\theta,t)\|_{\infty} \leq \|\omega\|_{\infty} + |K|\,\Gmax$. In particular, under \textup{A1--A3}, the GCC dynamics has a unique global strong solution on $[0,T]$ for every initial condition $\theta(0)$.
\end{theorem}

\begin{proof}[Proof sketch]
For each component, using $|\sin(x) - \sin(y)| \leq |x - y|$:
\begin{align}
|b_i(\theta,t) - b_i(\psi,t)|
&\leq |K|\,\Gmax \sum_j W_{ij}(t)\,\bigl|(\theta_j - \theta_i) - (\psi_j - \psi_i)\bigr| \notag\\
&\leq |K|\,\Gmax\,\Bigl[\sum_j W_{ij}(t)\,|\theta_j - \psi_j| + |\theta_i - \psi_i|\underbrace{\sum_j W_{ij}(t)}_{=\,1}\Bigr] \notag\\
&\leq 2\,|K|\,\Gmax\,\|\theta - \psi\|_{\infty}. \label{eq:lip_detail}
\end{align}
Standard results for SDEs with globally Lipschitz drift and bounded diffusion yield existence and uniqueness.
\end{proof}


\section{Finite-Horizon Stability}

\begin{theorem}[Stability under drift perturbations]
\label{satz:stabilitaet}
Consider two systems driven by the same Brownian motion:
\begin{align}
\mathrm{d}\theta(t) &= b\bigl(\theta(t),t\bigr)\,\mathrm{d}t + \Sigma(t)\,\mathrm{d}B(t), \\
\mathrm{d}\vartheta(t) &= \tilde{b}\bigl(\vartheta(t),t\bigr)\,\mathrm{d}t + \Sigma(t)\,\mathrm{d}B(t).
\end{align}
Suppose both drifts are globally $L$-Lipschitz and
\begin{equation}
\|b(x,t) - \tilde{b}(x,t)\|_{\infty} \leq \varepsilon(t) \quad \text{for all } x \in \mathbb{R}^N,\; t \in [0,T].
\end{equation}
Then
\begin{equation}
\label{eq:gronwall}
\boxed{\;
\|\theta(t) - \vartheta(t)\|_{\infty}
\;\leq\; e^{Lt}\,\biggl(\|\theta(0) - \vartheta(0)\|_{\infty} + \int_0^t \varepsilon(s)\,\mathrm{d}s\biggr)
\;}
\end{equation}
\end{theorem}

\begin{proof}[Proof sketch]
The diffusion terms cancel in the difference equation. For $\delta(t) = \theta(t) - \vartheta(t)$, we have $\|\delta'(t)\|_{\infty} \leq L\,\|\delta(t)\|_{\infty} + \varepsilon(t)$. Gronwall's inequality yields the claim.
\end{proof}

\begin{significance}
Small changes in $W$, $G$, $g_i$, or $\varphi_{ij}$ produce only controlled changes in the trajectory on finite time horizons.
\end{significance}


\section{Continuity of the Observables}

\begin{proposition}[Local Lipschitz property]
\label{prop:stetigkeit}
For fixed $\tau_D$, $\tau_M$, $\lambda$, and $\varepsilon_{\mathrm{floor}}$, the observables $R(t)$, $\Deff^{\lambda}(t)$, and $\Mtau(t)$ are continuous functions of the window trajectory. On compact sets they are locally Lipschitz. In particular, on any suitable compact tube there exists a constant $\Lobs$ such that
\begin{equation}
\label{eq:Lobs}
\|O(\mathrm{traj}_1) - O(\mathrm{traj}_2)\|_{\infty} \;\leq\; \Lobs \cdot \|\mathrm{traj}_1 - \mathrm{traj}_2\|_{\infty},
\end{equation}
where $O = (R,\;\Deff^{\lambda},\;\Mtau)$.
\end{proposition}

\paragraph{Empirical estimation of $\Lobs$.}
One chooses small initial perturbations $\delta_0$, simulates $\theta$ and $\theta^{\delta}$ with the same noise path, and computes
\begin{equation}
\hat{L}_{\mathrm{obs}}(\delta_0) = \frac{\sup_t \|O(\theta^{\delta})(t) - O(\theta)(t)\|_{\infty}}{\sup_t \|\theta^{\delta}(t) - \theta(t)\|_{\infty}}.
\end{equation}


\section{Residual-Backbone Re-Entry}

\begin{theorem}[Sufficient condition for re-entry]
\label{satz:reentry}
Let $S = S(t_1;\,g_*)$ be the residual backbone structure defined in the main text. Consider, on $I = [t_1,t_2]$, the system restricted to $S$ in which the coupling terms coming from outside $S$ are switched off. Define the external residual mass
\begin{equation}
\alpha_S(t) = \max_{i \in S}\sum_{j \notin S} W_{ij}(t).
\end{equation}
Assume that:
\begin{enumerate}
\item The reduced system lies on $I$ within the access region with safety margin $\rho > 0$:
\begin{align}
R_S(t) &> \Rmin + \rho, \notag\\
\Dmin + \rho &< D_{\mathrm{eff},S}^{\lambda}(t) < \Dmax - \rho, \label{eq:marge}\\
\Mmin + \rho &< M_{\tau,S}(t) < \Mmax - \rho. \notag
\end{align}
\item The initial states coincide at $t_1$.
\item A local Lipschitz constant $\Lobs$ for $O = (R_S,\,D_{\mathrm{eff},S}^{\lambda},\,M_{\tau,S})$ exists.
\item The following holds:
\begin{equation}
\label{eq:eta}
\boxed{\;
\int_{t_1}^{t_2} \alpha_S(s)\,\mathrm{d}s \;<\; \eta, \qquad
\eta \;<\; \frac{\rho}{\Lobs \cdot |K| \cdot \Gmax \cdot e^{2|K|\Gmax(t_2 - t_1)}}
\;}
\end{equation}
\end{enumerate}
Then the full system restricted to $S$ remains in the access region throughout $I$.
\end{theorem}

\begin{proof}[Proof sketch]
By Theorem~\ref{satz:stabilitaet} with $L = 2|K|\Gmax$ and $\varepsilon(t) = |K|\Gmax\,\alpha_S(t)$,
\begin{equation}
\sup_{t \in I}\|\theta_{\mathrm{full}}(t) - \theta_{\mathrm{red}}(t)\|_{\infty}
\;\leq\; |K|\,\Gmax\,e^{2|K|\Gmax(t_2 - t_1)}\cdot\eta.
\end{equation}
Using $\Lobs$, it follows that $\|O_{\mathrm{full}} - O_{\mathrm{red}}\|_{\infty} \leq \Lobs \cdot |K|\,\Gmax\,e^{2|K|\Gmax(t_2-t_1)}\cdot\eta$. If $\eta$ is chosen according to~\eqref{eq:eta}, this stays below $\rho$.
\end{proof}

\begin{significance}
This is the mathematically sharpest form of the selective-preservation re-entry extension: a transient re-entry is sufficiently possible whenever a residual subnetwork can carry the dynamics and the rest of the system does not perturb it too strongly on a short interval.
\end{significance}

\paragraph{Conservativity.}
The $\eta$-condition is a Gronwall-based worst-case bound. If numerical dynamics show a re-entry even though the analytical bound is not satisfied, this primarily reflects conservativity of the bound, not evidence against the simulated mechanism.


\section{Corollary: Gain-Mediated Lucidity-Like Re-Entry}

\begin{corollary}
\label{kor:lucidity}
If there exist a subset $S$ and an interval $I$ such that $\alpha_S(t)$ becomes small on $I$, the effective drives $\Lambda_i(t) = K\,G(t)\,g_i(t)$ for $i \in S$ do not fall below the reduced access threshold, and the reduced dynamics on $S$ remains in the access region with safety margin, then a transient re-entry of the full system onto $S$ is internally sufficiently possible within the model.
\end{corollary}

The mathematical core of the hypothesis:

\begin{quote}
\itshape
Lucidity-like re-entry, if generated by the GCC mechanism, can be understood as finite-horizon persistence of a residual backbone regime under decreasing external perturbation load.
\end{quote}


\section{What This Appendix Delivers and What It Does Not}

\paragraph{Delivers.}
Well-posedness of the GCC as an SDE; controlled bounds for the core observables; finite-horizon stability; a genuine sufficient re-entry condition (Theorem~\ref{satz:reentry}).

\paragraph{Does not deliver.}
No mathematical proof that real terminal lucidity exists; no general theorem that every form of dementia shifts the access region; no proof of a universal threshold $K_c$; no identifiability result for real EEG/MEG data.


% ==================================================================
\end{document}
